% cabeçalho genérico:

\usepackage[brazil]{babel}
\usepackage[utf8]{inputenc}
\usepackage[top=3cm, bottom=3.5cm, left=3cm, right=2cm]{geometry}
\usepackage[nodisplayskipstretch]{setspace}
\usepackage[T1]{fontenc}
\usepackage{tikz}

\usepackage{graphicx,xcolor,comment,enumerate,multirow,multicol} 

\usepackage{amsmath,amsthm,amsfonts,amssymb,dsfont,mathtools}

\usepackage{hyperref}	
\hypersetup{
	colorlinks=true,
	linkcolor=cyan,
	filecolor=cyan,      
	urlcolor=cyan,
}

\usepackage{enumitem, float}

\graphicspath{{imagens}}

\usepackage{parskip}

\setlength{\parindent}{1.0cm}
\setlength{\parskip}{2pt}


% Inclui o pacote xcolor com a opção para nomes de cores
\usepackage[svgnames]{xcolor}

% Define a cor do texto usando o código hexadecimal
\definecolor{Cornsilk}{HTML}{FFF8DC}

% Define a cor de fundo da página como preta
\pagecolor{Black}

% Define a cor padrão do texto para todo o documento
\color{Cornsilk}